% TOP_PROBLEM: {"id": 30003897, "title": "Green’s Canonical Syzygy Conjecture for Arbitrary Curves", "category_id": 5, "category": {"id": 5, "name": "algebraic_geometry", "display_name": "Algebraic Geometry", "description": "Geometric objects defined by polynomial equations.", "slug": "algebraic-geometry", "order_index": 5, "created_at": "2026-07-31T15:26:25.671Z"}, "status": "open"}
% FIRST_POSTED: 2026-09-27
% !TeX program = lualatex
% ATTEMPT_STATUS: unresolved
\documentclass[11pt]{article}
\usepackage[margin=25mm]{geometry}
\usepackage{fontspec}
\setmainfont{Latin Modern Roman}
\usepackage{amsmath,amssymb,mathtools}
\usepackage{unicode-math}
\setmathfont{Latin Modern Math}
\usepackage{xurl,hyperref}
\hypersetup{colorlinks=true,urlcolor=blue}
\setcounter{secnumdepth}{0}
\setlength{\parindent}{0pt}
\setlength{\parskip}{5pt}
\setlength{\emergencystretch}{3em}
\begin{document}
\section*{\texorpdfstring{Green’s Canonical Syzygy Conjecture for Arbitrary Curves}{Green’s Canonical Syzygy Conjecture for Arbitrary Curves}}\label{30003897}
\subsection{Definitions and mathematical statement}
In the complex characteristic-zero setting of the cited source, let $C$ be a smooth nonhyperelliptic curve of genus $g$, $V=H^0(C,K_C)$, and define $K_{p,1}(C,K_C)$ as the middle cohomology of the Koszul complex
\[
 \bigwedge^{p+1}V\otimes H^0(C,\mathcal O_C)
 \longrightarrow\bigwedge^pV\otimes H^0(C,K_C)
 \longrightarrow\bigwedge^{p-1}V\otimes H^0(C,K_C^2).
\]
The maps alternate multiplication of sections. The Clifford index is the minimum of $\deg L-2(h^0(L)-1)$ over line bundles with $h^0(L),h^1(L)\ge2$, with the standard value one for a nonhyperelliptic genus-three curve. The selected Green equivalence is
\[
 [K_{\ell,1}(C,K_C)=0\text{ for every }\ell\ge p]
 \quad\Longleftrightarrow\quad\operatorname{Cliff}(C)>g-p-2.
\]
The catalogue does not name a ground field; the source setting is stated explicitly rather than extending this assertion silently to arbitrary positive characteristic.
\par\smallskip\textit{Formulation and concept references:} \href{https://webusers.imj-prg.fr/~claire.voisin/Articlesweb/syzod.pdf}{[S1]}.

\subsection{Short English statement}
For every smooth nonhyperelliptic complex curve, does its Clifford index determine exactly where the linear syzygies of its canonical embedding vanish?
\subsection{Research attempt: canonical complete intersections and an index-boundary diagnostic}
\textbf{Status and index convention.} The all-curve, characteristic-zero Green target is \textbf{UNRESOLVED} by this attempt. We compute the complete linear strands for several canonical complete intersections and isolate the rank condition needed in general. There is also a boundary issue in the displayed catalogue formula: it does not state a range for $p$. If $p=0$ is included, a smooth plane quartic with the stipulated Clifford index one makes that literal equivalence false. We prove this explicitly below. The substantive research target is treated with positive integer indices $p\ge1$, as appropriate to the linear syzygy statement; the boundary diagnostic is not a disproof of the usual Green conjecture.

\textbf{Source audit.} The introduction of Voisin's primary paper [S1] prints the selected equivalence, identifies the Green--Lazarsfeld nonvanishing direction, and explains that its new result is for generic curves of odd genus. Its K3-surface construction and genericity cannot be erased to obtain the assertion for every complex curve. We use neither generic vanishing as a specialization theorem nor positive-characteristic examples as counterexamples to this characteristic-zero question. The computations below concern explicit geometric classes and are standard low-genus theory, with their algebra displayed in full.

\subsubsection{1. Reading the linear strand from a free resolution}
Let $S=\operatorname{Sym}V$, with $V$ in degree one, and
$R=\bigoplus_{q\ge0}H^0(C,K_C^q)$. The standard Koszul resolution of the residue field $\mathbb C$ over $S$ identifies
\[
 K_{p,q}(C,K_C)=\operatorname{Tor}^S_p(R,\mathbb C)_{p+q}.
                                                               \tag{1}
\]
Indeed its internal-degree $p+q$ part has consecutive terms
$\bigwedge^{p+1}V\otimes R_{q-1}$,
$\bigwedge^pV\otimes R_q$, and
$\bigwedge^{p-1}V\otimes R_{q+1}$, with the multiplication differentials in the question.

In a minimal graded free resolution of $R$, all matrix entries of the differentials lie in the homogeneous maximal ideal of $S$. Tensoring with $\mathbb C$ kills them. Thus a summand $S(-j)$ in homological degree $p$ contributes one dimension to $K_{p,j-p}$, and no contribution to a different internal degree. In particular the linear strand uses exactly the shifts $j=p+1$.

For a regular sequence of homogeneous polynomials $f_1,\ldots,f_c$ of degrees $a_1,\ldots,a_c$, the Koszul complex gives a minimal resolution of $S/(f_1,\ldots,f_c)$. Its term in homological degree $p$ is
\[
 \bigoplus_{|I|=p}S\Bigl(-\sum_{i\in I}a_i\Bigr).       \tag{2}
\]
Exactness follows inductively from the nonzerodivisor condition, by the mapping-cone argument adjoining one polynomial at a time. Minimality follows because each $a_i>0$. Formula (2), rather than a numerical guess from a Hilbert polynomial, will determine the syzygies below.

For the smooth projective complete intersections used here, their section rings agree with these homogeneous coordinate rings. One checks this by the sheafified Koszul resolution and the vanishing of intermediate cohomology of line bundles on projective space. This projective-space cohomology calculation is an external standard input; it is essential when translating a coordinate-ring resolution into the canonical Koszul complex of sections.

\subsubsection{2. Plane quartics and the omitted zero index}
Let $C\subset\mathbb P^2$ be a smooth plane quartic. Adjunction gives $K_C\cong\mathcal O_C(1)$, and $\deg K_C=4=2g-2$, so $g=3$. The canonical map is this plane embedding, hence $C$ is nonhyperelliptic. Its canonical ring is
\[
 R=\mathbb C[x_0,x_1,x_2]/(F_4),\qquad
 0\to S(-4)\xrightarrow{F_4}S\to R\to0.               \tag{3}
\]
The only nontrivial positive homological-degree Tor group occurs at $(p,j)=(1,4)$, that is, at $K_{1,3}$. Therefore
\[
 K_{p,1}(C,K_C)=0\qquad(p\ge1).                       \tag{4}
\]
The zero-index group also vanishes: its defining incoming map
$V\otimes H^0(\mathcal O_C)\to H^0(K_C)$ is the identity after choosing $1\in H^0(\mathcal O_C)$, so $K_{0,1}=0$.

By the genus-three convention in the root statement, $\operatorname{Cliff}(C)=1$. At $p=0$, the left-hand side of the displayed equivalence is true by (4) and $K_{0,1}=0$, while the right-hand side is $1>3-0-2=1$, which is false. This is a complete verification of the indexing defect if a universal quantifier over nonnegative $p$ was intended.

For every $p\ge1$, however, the same example satisfies the intended equivalence: all the groups on the left vanish, and $1>1-p$ is true. The minimal repair for this example is to specify the positive index range. Nothing in (3) addresses arbitrary genus, and it would be misleading to promote this small-index convention issue into a solution of the central conjecture.

\subsubsection{3. A genus-four family: one quadratic generator and no later linear syzygy}
Take a smooth curve of bidegree $(3,3)$ on a smooth quadric
$Q\cong\mathbb P^1\times\mathbb P^1\subset\mathbb P^3$.
The restriction of $\mathcal O_{\mathbb P^3}(1)$ to $Q$ is $\mathcal O_Q(1,1)$. Adjunction on $Q$ gives
\[
 K_C\cong\mathcal O_Q(1,1)|_C,\qquad
 \deg K_C=(1,1)\cdot(3,3)=6,
\]
so $g=4$ and the given embedding is canonical. The curve is cut out in $\mathbb P^3$ by the quadric equation $Q_2$ and a cubic $F_3$; restriction of ambient cubics onto $Q$ is surjective by the hypersurface sequence. These form a regular sequence. Hence
\[
0\to S(-5)\to S(-2)\oplus S(-3)\to S\to R\to0.       \tag{5}
\]
Equation (1) now gives
\[
 \dim K_{1,1}=1,\qquad K_{p,1}=0\quad(p\ge2).         \tag{6}
\]
The two projections $Q\to\mathbb P^1$ restrict to degree-three maps on $C$. Their line bundles have degree three and at least two sections; Riemann--Roch gives at least two sections for the residual canonical bundle as well. Thus $\operatorname{Cliff}(C)\le1$. Clifford's theorem, with its equality characterization, says that a smooth nonhyperelliptic curve has Clifford index at least one when the admissible set is nonempty. Therefore it is exactly one here.

For $p=1$, the left side of Green's equivalence fails because $K_{1,1}\ne0$, while $1>4-1-2=1$ is false. For $p\ge2$, (6) gives the left side and $1>2-p$ gives the right side. This verifies the entire positive-index equivalence for this genus-four family.

The cubic generator in (5) belongs to $K_{1,2}$, not $K_{1,1}$, and the last relation belongs to $K_{2,3}$, not $K_{2,1}$. Keeping internal degree separate from homological degree prevents an off-by-one inference about the linear strand.

\subsubsection{4. A genus-five family with Clifford index two}
Let $C\subset\mathbb P^4$ be a smooth complete intersection of three quadrics. Adjunction gives $K_C=\mathcal O_C(1)$; its degree is $2^3=8$, hence $g=5$. Its minimal resolution is
\[
0\to S(-6)\to S(-4)^{\oplus3}\to S(-2)^{\oplus3}
 \to S\to R\to0.                                    \tag{7}
\]
Consequently
\[
 \dim K_{1,1}=3,\qquad K_{p,1}=0\quad(p\ge2).         \tag{8}
\]
For a concrete family in which the Clifford index can also be checked geometrically, require one quadric to have rank four, for example
$Q=x_0x_1-x_2x_3$, and take two sufficiently general further quadrics avoiding its vertex $[0:0:0:0:1]$. Bertini on the smooth locus of $Q$ gives a smooth curve; the vertex is absent by construction.

The rank-four quadric is a cone over a smooth quadric surface. Either ruling of that surface gives a family of projective planes through the vertex, parameterized by $\mathbb P^1$. Away from the vertex this defines a morphism to $\mathbb P^1$. Its restriction to $C$ has degree four: a ruling plane meets the two remaining quadrics in a degree-four zero-dimensional complete intersection for a general plane. Thus $C$ has a basepoint-free $g^1_4$, and its residual bundle has $h^1=h^0\ge2$ by Riemann--Roch at degree four and genus five. This yields $\operatorname{Cliff}(C)\le2$.

There is no degree-three pencil on any smooth complete intersection of three quadrics. If such a pencil existed, a general fiber $D$ would consist of three distinct points, because the characteristic is zero. Riemann--Roch gives
$h^0(K_C-D)=h^0(D)+1\ge3$. Since $h^0(K_C)=5$, those three canonical points span at most a line: their evaluation map has rank at most two. They are distinct, so they lie on a projective line. Each quadric defining $C$ vanishes at three distinct points on that line and hence vanishes on the entire line. The line would then be contained in their scheme-theoretic intersection $C$, impossible for a smooth connected irreducible curve of degree eight. Connectedness and irreducibility here are the standard connectedness property of smooth positive-dimensional projective complete intersections.

To see that this excludes Clifford index one without an unstated classification, use residual duality. A line bundle computing index one can be replaced by $K_C\otimes L^{-1}$ without changing its Clifford value, and one of these two bundles has degree at most four. The equation
$\deg L-2(h^0(L)-1)=1$, together with $h^0(L)\ge2$, then forces $\deg L=3$ and $h^0(L)=2$. Removing a base point would give a degree-two pencil, contrary to the canonical embedding. Thus an index-one bundle would produce the forbidden degree-three pencil. Nonhyperellipticity excludes index zero by Clifford's theorem. Hence the chosen family has Clifford index exactly two.

At $p=1$, $K_{1,1}\ne0$ and $2>5-1-2=2$ is false. At every $p\ge2$, (8) holds and $2>3-p$ is true. Thus this family also verifies every positive-index case of the displayed equivalence.

\subsubsection{5. Hilbert-series checks do not determine arbitrary linear strands}
For a canonical curve of genus $g$, Riemann--Roch gives
\[
 h^0(K_C^q)=(2q-1)(g-1)\quad(q\ge2),\qquad
 h^0(K_C)=g.
\]
The Hilbert series is therefore
\[
 H_R(t)=\frac{1+(g-2)t+(g-2)t^2+t^3}{(1-t)^2}.         \tag{9}
\]
For the three complete intersections, the resolution numerators are respectively
$1-t^4$, $(1-t^2)(1-t^3)$ and $(1-t^2)^3$, over denominators $(1-t)^3,(1-t)^4,(1-t)^5$. Cancelling the appropriate factors of $1-t$ reproduces (9) for $g=3,4,5$. Exact polynomial multiplication is a useful independent diagnostic of every shift in (3), (5), and (7).

For an arbitrary canonical curve, however, a Hilbert series gives only alternating sums of graded Betti numbers:
\[
 (1-t)^gH_R(t)=\sum_{p,j}(-1)^p\beta_{p,j}t^j.
\]
Contributions in adjacent homological degrees can cancel at the same internal degree. Knowing this polynomial therefore does not recover every $\beta_{p,p+1}$. The complete-intersection examples avoid that ambiguity because their whole minimal resolutions were constructed, not merely inferred from (9).

The same distinction arises if one knows only the number of quadrics. Even a complete count of $K_{1,1}$ does not force all higher $K_{p,1}$ to vanish: relations among those quadrics are additional data. A proof for general curves must control the actual multiplication differentials or geometric structures that govern them.

\subsubsection{6. The rank formulation and the genericity obstruction}
For fixed $p$, write the two differentials in the question as $A_p$ and $B_p$. Since $B_pA_p=0$,
\[
 \dim K_{p,1}
 =\dim\bigl(\textstyle\bigwedge^pV\otimes V\bigr)
   -\operatorname{rank}A_p-\operatorname{rank}B_p.      \tag{10}
\]
Thus vanishing is a maximal-rank equality subject to the Koszul relation. In a family with the relevant spaces organized as vector bundles, a rank may decrease on a closed determinantal locus. Vanishing at a generic parameter is then an open condition, not a conclusion at every special parameter.

The one-by-one matrix $[t]$ already demonstrates the direction: its kernel is zero for $t\ne0$ and one-dimensional at $t=0$. This is not a family of canonical curves, but it exactly refutes the purely linear-algebra inference that generic maximal rank persists under all specializations. In the geometric problem, special curves also have special linear series and possibly a different Clifford index, so the desired theorem must compare the two jumps in a controlled way.

A possible successful route would establish the needed ranks uniformly on each Clifford-index stratum, or show that every rank-drop locus at the forbidden indices is accounted for by a lower-Clifford linear series. The primary source proves powerful generic cases, but the adjective generic cannot be removed by continuity alone. No such uniform rank argument is supplied here.

\subsubsection{7. Verification and final scope}
All resolutions used above come from regular sequences and have explicitly positive-degree entries, ensuring minimality. Adjunction checks that the embedding line bundle is actually canonical. The genus-five ruling pencil avoids the cone vertex, and the no-trigonal argument uses a general reduced fiber rather than assuming that every degree-three divisor is reduced. Hilbert numerators and the resulting $K_{p,1}$ tables are checked exactly in the saved diagnostic.

For the literal nonnegative-index reading, the plane-quartic calculation disproves the $p=0$ extension with the stated genus-three convention. For the repaired positive-index target, the complete-intersection subcases are proved and the general problem remains open in this attempt. The first missing step is a uniform rank or nonvanishing-to-linear-series theorem covering every relevant special curve, not just a generic member of moduli. The substantive Green target is \textbf{UNRESOLVED} here.

\subsection{Sources}
\begin{itemize}
\item[S1]Claire Voisin, Green's canonical syzygy conjecture for generic curves of odd genus, Compositio Mathematica, 2005.
\url{https://webusers.imj-prg.fr/~claire.voisin/Articlesweb/syzod.pdf}.
\textit{Formulation source.}
\end{itemize}
\end{document}
